% !PRINTED 131
\weberpage{131}

If the inequalities (13) are written in the form
\[
\begin{aligned}
 a-a'&>\frac{a^2}{4Q},\\
 a'-a''&>\frac{a'^2}{4Q},\\
 &\hspace{1.5em}\vdots\\
 a^{(\nu-1)}-a^{(\nu)}&>\frac{(a^{(\nu-1)})^2}{4Q},
\end{aligned}
\]
and added from the $(\mu+1)$-st to the $(\nu+1)$-st, one obtains
\[
 (a^{(\mu)})^2+(a^{(\mu+1)})^2+\cdots+(a^{(\nu)})^2
 <4Q\bigl(a^{(\mu)}-a^{(\nu+1)}\bigr),
\]
and therefore, when $\mu$ and $\nu$ both tend to infinity,
\begin{equation}
 \lim\bigl((a^{(\mu)})^2+(a^{(\mu+1)})^2+\cdots+(a^{(\nu)})^2\bigr)=0.
\tag{16}
\end{equation}

For the sequence of numbers $\alpha,\alpha',\alpha'',\ldots$ the law of formation resulting from (11) is
\begin{equation}
\begin{aligned}
 \alpha-\alpha'&=\frac{a f(\alpha)}{2Q f'(\alpha)},\\
 \alpha'-\alpha''&=\frac{a' f(\alpha')}{2Q f'(\alpha')},\\
 &\hspace{1.5em}\vdots\\
 \alpha^{(\nu-1)}-\alpha^{(\nu)}&=
 \frac{a^{(\nu-1)} f(\alpha^{(\nu-1)})}{2Q f'(\alpha^{(\nu-1)})}.
\end{aligned}
\tag{17}
\end{equation}
If an arbitrary number of these equations are added, from the $(\mu+1)$-st to the $(\nu+1)$-st, then
\begin{equation}
 \alpha^{(\mu)}-\alpha^{(\nu+1)}=
 \frac{a^{(\mu)}f(\alpha^{(\mu)})}{2Q f'(\alpha^{(\mu)})}+\cdots+
 \frac{a^{(\nu)}f(\alpha^{(\nu)})}{2Q f'(\alpha^{(\nu)})}.
\tag{18}
\end{equation}
Since the inequalities
\[
 |f'(\alpha)|>k
\]
are satisfied, it follows for the absolute value of this difference that
\begin{equation}
 |\alpha^{(\mu)}-\alpha^{(\nu+1)}|<
 \frac{(a^{(\mu)})^2+(a^{(\mu+1)})^2+\cdots+(a^{(\nu)})^2}{2Qk},
\tag{19}
\end{equation}
and hence, by (16),
\begin{equation}
 \lim |\alpha^{(\mu)}-\alpha^{(\nu+1)}|=0.
\tag{20}
\end{equation}

This says (according to the concept of a sequence of numbers explained in the Introduction, p. 15) that the numbers $\alpha,\alpha',\alpha'',\ldots$ approach a definite limit, which we shall denote by $\xi$.
%CONT>>
